% Section VI: Phase 2 Gap and Conclusion. Numbers: results/metrics/
% test_results.csv, feature_check.csv, error_breakdown.csv, per_phone.csv and
% cv_check.csv (all repeated from Section V).
\section{Phase 2 Gap and Conclusion}\label{sec:conclusion}

\emph{The gap.} Our three task models are separate and share nothing, and the
results point to what this misses. First, no model is best at every task:
XGBoost finds the most floors (95.2\%), Random Forest the best positions
(7.42~m). Second, AP-affinity, a
hand-made way of sharing location knowledge (where each WAP is heard) with all
three tasks, gave our largest gain (CV error from 15.87 to 11.38~m); a model
that learns such shared knowledge by itself may gain more. Third, most of the
remaining error is position error on the right floor (86.2\% of the Random
Forest's total error), and new phones are harder than seen ones (8.30~m
against 5.84~m).

\emph{Phase~2} will train one multi-task neural
network~\cite{caruana1997multitask} that learns building, floor and position
together: shared layers used by all three tasks, then three output heads. Kim
et~al.~\cite{kim2018scalable} showed that one network with building, floor and
location outputs can be trained on UJIIndoorLoc. We will evaluate our network
with the same user-group CV, test set and metrics, and compare it with our
Phase~1 models.

\emph{Conclusion.} Fair testing matters most: random CV promised 3.67~m, but
the Baseline RF reached 12.52~m on the test set. Chosen by user-group CV, the
final Random Forest with AP-affinity found every building and 92.2\% of the
floors, with a mean 2D error of 7.42~m.
